\chapter*{Preface}
\addcontentsline{toc}{chapter}{Preface}

\noindent
Every society produces goods, laws, and infrastructures.
Every society also produces beliefs---about what people are, what they deserve, what is possible, and what is natural.
The first kind of production is studied by economics, engineering, and political science.
The second kind has no dedicated science.

This book proposes one.

\textbf{Doxonomics}---from the Greek \textit{doxa} (opinion, common belief, public orthodoxy) and \textit{-nomics} (systematic study of allocation and structure)---is the science of how material power shapes public belief.
Its central question is not whether a belief is true, but what it cost to make it appear self-evident: who financed it, which institutions processed it, through what channels it circulated, and what social effects followed from its dominance.

The need for such a field is not difficult to motivate.
Global advertising exceeded one trillion dollars in 2024.
Lobbying, think-tank funding, public relations, and curriculum sponsorship add hundreds of billions more.
Existing disciplines study parts of this process: sociology of knowledge examines the social formation of ideas, political communication traces agenda-setting, media studies analyzes representation, propaganda studies maps persuasion techniques, and social epistemology investigates the conditions of collective knowledge.
What no existing field does is integrate these strands around \emph{explicit accounting of material support, institutional routing, and measurable belief effects}.
Political economy studies the distribution of material power but not its conversion into common sense.
Sociology of knowledge studies the social formation of ideas but rarely performs financial accounting on them.
Media studies analyzes representation but seldom asks which representations are continuously subsidized and to what measurable effect.
Doxonomics sits at the intersection of all these fields and adds what they individually lack: a unified framework for auditing the industrial production of belief.

This textbook is mathematically serious.
We take the position that critique without formal tools remains vulnerable to the same vagueness it diagnoses.
The field needs stock-and-flow models, causal inference, network analysis, information-theoretic measures, and---where the structure demands it---sheaf-theoretic coherence conditions.
Many of the models presented here are heuristic or conjectural rather than empirically calibrated; we mark them as such.
Their purpose is to clarify structure and generate testable predictions, not to simulate maturity the field has not yet earned.
Where a result is a proven theorem, we say so.
Where it is a stylized model, a conjecture, or a candidate formalization, we say that too.

Throughout the text, we reference interactive \emph{playgrounds}: computational environments where the reader can manipulate funding parameters, observe belief diffusion, simulate narrative cascades, and test the sensitivity of doxonomic models to their assumptions.
These are available at the Piatra Institute's companion site.

The book proceeds in five parts.
Part~I lays the conceptual and historical foundations.
Part~II develops the core formal apparatus---epistemic capital, narrative infrastructure, common-sense capture, doxonomic power, and sheaf-like coherence.
Part~III introduces the methods: measurement, accounting, network analysis, causal inference, computational text analysis, and experimental design.
Part~IV applies the machinery to nine domains, from the managed narrative of selfish human nature to crisis doxonomics.
Part~V turns to resistance, repair, and the ethics of studying belief production without becoming belief engineers.

No single author could write this book.
It emerges from a research program at the Piatra Institute, and it will remain open to revision as the field develops.
The reader is invited not only to learn doxonomics but to extend it.

\vspace{1em}
\noindent\textit{Piatra Institute, \the\year}
